% CoRL 2026 Workshop PhysWM, 4 pages excluding references, double anonymous, non-archival.
% The workshop encourages the CoRL template; its 2026 style file is not publicly
% downloadable, so this uses a neutral single-column layout of the same family.
\documentclass[10pt,letterpaper]{article}

\usepackage[letterpaper,top=1in,bottom=1in,left=1.25in,right=1.25in]{geometry}
\usepackage{times}
\usepackage[numbers,sort&compress]{natbib}
\usepackage{graphicx}
\usepackage{amsmath,amssymb}
\usepackage{booktabs}
\usepackage[hidelinks]{hyperref}
\usepackage[capitalize]{cleveref}
\usepackage{titlesec}
\titlespacing*{\section}{0pt}{1.1ex plus .2ex}{0.7ex}
\titlespacing*{\paragraph}{0pt}{0.8ex}{1em}
\setlength{\parskip}{0pt}
\providecommand{\etal}{et al.}
\newcommand{\sys}{the engine}

\title{\vspace{-3.5ex}Better Physics, Same Score\vspace{-1ex}}
\author{Anonymous submission\vspace{-2ex}}
\date{}

\begin{document}
\maketitle

\begin{abstract}
\noindent
A world model is useful for control only if it gets the physics right, so we need benchmarks
that reward physical accuracy rather than plausible-looking output. We test whether the most
widely used counterfactual video benchmark does. We build an engine that answers CLEVRER's
counterfactual questions by intervention, with the simulator as a swappable component, and
run 13 simulators through it on 1,398 held-out videos. A 2.2$\times$ spread in rollout error
moves the benchmark score by 2.2 points, straight-line motion scores within 0.75 points of a
learned world model, and on predictive questions the ordering reverses: the most accurate
simulator scores lowest. What does move the score is evaluation logic that contains no
physics at all, worth up to 4.6 points. Physics as a whole matters, since removing simulation
costs 9.0 points, but accuracy above a low bar buys almost nothing. We argue that a benchmark
meant to verify a world model must report, next to accuracy, the share of decisions the
simulator actually made and the share whose evidence is valid, and we show both are
measurable.
\end{abstract}

\vspace{-1ex}
\section{The problem this measures}
\vspace{-0.5ex}

This workshop asks what benchmarks can verify that a world model is control-useful rather
than only visually convincing. The question presumes that current benchmarks reward physical
accuracy. We tested that presumption on CLEVRER~\cite{yi2020clevrer}, the most used benchmark
for counterfactual physical reasoning, and found that its counterfactual score is nearly
blind to the accuracy of the simulator that produces the answers.

The result matters beyond one dataset. Counterfactual and predictive questions reduce a
continuous trajectory to a discrete outcome: did these two objects touch? A yes-or-no outcome
changes only when a trajectory error moves a pair across the contact distance. Everywhere
else, halving the error changes nothing that is scored. Manipulation benchmarks that grade
task success share this structure, so the measurement we report is a property of outcome-
scored evaluation, not of CLEVRER alone.

\vspace{-0.5ex}
\section{An engine with the simulator as a variable}
\vspace{-0.5ex}

To attribute a score to the simulator, everything else must be held fixed. Our engine answers
a counterfactual with the three steps of an intervention~\cite{pearl2009causality}: the
recording plays the role of abduction, the edit is the action, and prediction is a
simulation. Every object follows its recorded path until the intervention can reach it, then
is handed to a simulator a few frames before it would lose a recorded contact, and the world
runs on past the end of the video. The engine returns the event that decided each answer: the
two objects, the frame, and whether it was recorded or simulated. It is never trained on
questions or answers. Simulating only the objects an intervention affects follows
CRCG~\cite{ishay2024crcg}; building one reusable editable world per video is also the design
of PhysMind~\cite{yang2026physmind}.

The simulator is a module. We plug in a learned world model with a pairwise potential, a
symplectic integrator and learned per-object friction; textbook laws whose constants are
fitted on training videos, in the spirit of system identification~\cite{ding2021vrdp};
straight-line motion with and without friction; nine earlier checkpoints of the learned
model; and six retrained seeds. Rollout error is measured separately, on held-out validation
videos, as mean position error at 40 frames.

\vspace{-0.5ex}
\section{What the benchmark rewards}
\vspace{-0.5ex}

\begin{figure}[t]
\centering
\includegraphics[width=\linewidth]{figures/fig3_audit.pdf}
\vspace{-3.5ex}
\caption{Physics quality against benchmark score for 13 simulators. (a) Counterfactual
questions: a 2.2$\times$ spread in rollout error moves the score by 2.2 points, and most
differences sit inside the 95\% intervals. (b) Predictive questions: the ordering reverses
and the most accurate simulator scores lowest. Error on 200 validation videos, scores on
1,398 held-out test videos.}
\label{fig:audit}
\vspace{-2ex}
\end{figure}

\cref{fig:audit}a plots the counterfactual score of all 13 simulators against their rollout
error. The error spans 0.050 to 0.113, a factor of 2.2, while the score spans 89.2\% to
91.4\% of options (Pearson $r=-0.59$). Straight-line motion has 1.8 times the rollout error
of the learned model and scores 0.75 points lower. On predictive questions the relation
reverses (\cref{fig:audit}b): straight lines answer 91.4\% of options while the fitted laws,
which have the lowest rollout error, answer 88.1\% ($z=-3.8$). No simulator is best on both
question types.

The reason is visible in the answers. The four main simulators give the same answer on
91.4\% of the 9,350 counterfactual options, although at least one asked object is simulated
before the video ends on 47.8\% of them. Every difference between simulators comes from the
remaining 8.6\%, which suggests most asked pairs sit far from the contact boundary where a
trajectory error could flip an answer.

This is not an argument that physics is unnecessary. Removing the simulation entirely costs
9.0 points per option and 22.2 per question. The claim is narrower and, we think, more
useful: \emph{above a low bar, accuracy buys almost nothing on this metric}.

\paragraph{What does move the score.} Two choices that contain no physics matter more than
the choice of simulator. Stopping the simulation at the last recorded frame costs 4.6 points
per option and 14.3 per question, because CLEVRER's counterfactual answers include events
after the video ends. Handing objects to the simulator 12 frames before a lost contact
instead of 3 costs a further 1.2 and 3.2 points: every extra frame of real recorded path is
drift avoided. A system can therefore gain more from when it starts simulating than from how
well it simulates.

\begin{table}[t]
\centering\small
\setlength{\tabcolsep}{5pt}
\caption{Counterfactual accuracy on 1,398 held-out videos (9,350 options, 2,619 questions)
against rollout error measured separately. Sorted by rollout error.}
\label{tab:sims}
\begin{tabular}{@{}lccc@{}}
\toprule
Simulator & Rollout error @40 & Options (\%) & Questions (\%) \\
\midrule
textbook laws, fitted on training videos & 0.050 & \textbf{91.4} & \textbf{75.0} \\
learned world model & 0.063 & 90.4 & 72.3 \\
straight lines + friction & 0.095 & 90.2 & 71.9 \\
straight lines & 0.113 & 89.7 & 70.2 \\
\midrule
no simulation at all & -- & 81.4 & 50.1 \\
a 7B vision-language model, 16 frames & -- & 54.4 & 2.9 \\
answer ``no'' to every event & -- & 54.1 & 0.0 \\
\bottomrule
\end{tabular}
\vspace{-2.5ex}
\end{table}

\cref{tab:sims} places the spread in context. The gap between the best and worst real
simulator is 1.7 points; the gap between having a simulator and not having one is 9.0; and a
general vision-language model given the frames lands at the level of answering ``no'' to
everything.

\vspace{-0.5ex}
\section{Two numbers a benchmark should report}
\vspace{-0.5ex}

Because the engine returns the event behind each answer, we can ask what the accuracy rests
on. Of the correct answers whose deciding event can be checked against CLEVRER's annotation,
98.2\% rest on a real recorded event, meaning an annotated collision of the same pair within
15 frames or no collision of that pair at all. That is reassuring. The other half of the
story is not: 40.5\% of correct answers are decided in the simulated part of the world, where
the dataset has no ground truth, so they can be inspected but never verified.

We propose that benchmarks intended to verify world models report two quantities alongside
accuracy:

\vspace{-0.8ex}
\begin{itemize}\itemsep1pt \parskip0pt
\item \textbf{the share of decisions the simulator made}, separating answers that the
recording already settles from answers the physics produced; and
\item \textbf{the share of correct answers whose evidence is valid}, checked against
annotation where it exists.
\end{itemize}
\vspace{-0.8ex}

Both are cheap to compute once a system reports its deciding event, and together they
distinguish a model that is right for the right reason from one that inherits its score from
the evaluation protocol. On CLEVRER these two numbers reveal that a large, independently
verified improvement in physics, a 25\% reduction in rollout error from adding dissipation,
moves the benchmark by nothing measurable.

\paragraph{Applying this to a manipulation benchmark.} The recipe needs no new annotation.
Make the simulator a swappable component and run at least one deliberately weak baseline
through the same harness, for instance constant-velocity extrapolation or a rigid-body
simulator with default parameters. If task success barely separates them, the benchmark is
scoring something other than the dynamics, and the separation it does show tells you how much
of the headline number the physics earned. Then, for every scored episode, record which
predicted event the decision rested on and whether that event occurred in the real rollout.
The first measurement costs one extra run per baseline; the second costs a log line. Together
they turn \emph{our model scores higher} into \emph{our model scores higher because it
predicted contacts that actually happened}, which is the claim a control result needs.

\vspace{-0.5ex}
\section{Limitations}
\vspace{-0.5ex}

The scenes are synthetic, planar, with three shapes and one object size, and no manipulator
is involved; we have not tested the engine on real video or on contact-rich manipulation. Our
perception reads a detector trained with the dataset's own labels. The 13 simulators span a
range of quality but all are rigid-body planar models, so we cannot say where the saturation
begins, only that it is below the accuracy of straight-line motion. Whether the same
insensitivity holds for manipulation benchmarks that grade task success is the open question
we would most like to discuss at this workshop.

{\small
\bibliographystyle{plainnat}
\bibliography{refs}
}

\end{document}
